\chapter{The Finite-Context Theorems}\label{ch09:main}

This chapter proves the finite-context theory: the four registered theorems NT2-1 through NT2-4, together with the propriety proposition that separates the successor's generators from the legacy collapse. Every combinatorial input was computed in Chapter~7; every provenance question was quarantined into Chapter~8; what remains is the mathematics of sites and sheaves on the sixteen-object lattice $L=\mathcal{P}(\Omega)$, carried out at full length. The associated-sheaf reflectors and the composition theorem---which require constructing sheafification explicitly---are deferred to Chapter~10; the executable corroboration of everything is Chapter~11.

\section{Standing of the results and ledger discipline}\label{ch09:sec-standing}

Each theorem below is registered in the Noetic Topos 2.0 theorem ledger (\texttt{NOETIC\_TOPOS\_2\_THEOREMS\_v0.1.json}, schema \texttt{noetic\_topos2.theorem\_ledger.v0.1}, as of 2026-08-04, document status \textsc{draft / gate 1 research candidate / not licensed}). This monograph reproduces, for each theorem, the ledger record's identity and status fields verbatim: its \texttt{id}, statement kind, proof status, review state, novelty status, admission ring, and claim-boundary text. Two disciplines govern the reproduction.

First, \emph{status fidelity}. Every finite-context record carries proof status \texttt{INFORMAL\_PROOF\_PRESENT}, review state \texttt{IN\_REVIEW}, novelty status \texttt{NOVELTY\_UNASSESSED}, and admission ring 4. The proofs in this chapter are those informal proofs, expanded; writing them out at monograph length changes none of the statuses. In particular: these are not peer-reviewed theorems, they are not proof-assistant-verified beyond the bounded Lean scope described in Chapter~11, and no novelty in the scholarly sense is claimed for any of them. What entitles the word ``theorem'' here is the registered statement kind together with the finite, exhaustively checkable character of the claims; what forbids more is the ledger.

Second, \emph{corroboration is not verification}. Each theorem is corroborated by the executable finite certificate (\texttt{FINITE\_MODEL\_CERTIFICATE\_v0.1.json}, model \textsc{nt2-context-2x2-transverse-partitions}): a dependency-free checker exhausts all sixteen contexts, all sixteen representables, both witness presheaves, and every generating-cover pullback. Where a certificate check name is relevant we cite it. The certificate's own claim boundary governs every such citation: exhaustive execution over one sixteen-object finite context site ``is not proof-assistant verification, external peer review, empirical validation, or a product-claim license.''

\begin{notation}\label{ch09:not-recall}
Throughout, $L$, $\Omega$, the atoms $\omega_{ij}$, the partitions $A=\{A_{0},A_{1}\}$, $R=\{R_{0},R_{1}\}$, $Q$ (singletons), the generating families $K_{A}(V)$, $K_{R}(V)$, $K_{Q}(V)$ with retained empty pieces, and the null cover at $\varnothing$ are as in Chapter~7 (Definitions~\ref{ch07:def-omega}, \ref{ch07:def-lattice}, \ref{ch07:def-partitions}, \ref{ch07:def-generating}). Presheaves are ordinary small $\Set$-valued functors $L^{\mathrm{op}}\to\Set$ in a fixed universe. For a sieve $\sigma$ on $V$ (equivalently a down-closed family of subsets of $V$, Definition~\ref{ch07:def-sieve-poset}) and a presheaf $F$, a \emph{matching family} for $\sigma$ is a choice of $x_{S}\in F(S)$ for each $S\in\sigma$ with $x_{S}|_{T}=x_{T}$ whenever $T\subseteq S$; for a finite indexed family $\{S_{k}\to V\}$ a matching family is a choice of $x_{k}\in F(S_{k})$ agreeing on all pairwise intersections $S_{k}\cap S_{l}$. The \emph{gluing map} sends $x\in F(V)$ to its family of restrictions. $F$ satisfies the sheaf condition for the family (or sieve) when the gluing map is a bijection onto matching families \cite{maclane, stacks00ZC}.
\end{notation}

\section{The coverage-basis lemma}\label{ch09:sec-basis}

The passage from generating families to a Grothendieck topology runs through the standard notion of a coverage basis (pretopology): a system of families containing identity covers, stable under pullback, and stable under composition, whose generated sieve-collections then satisfy the topology axioms \cite{stacks00ZC, johnstone}. We verify each basis axiom for each partition, in all cases.

\begin{lemma}[Coverage basis]\label{ch09:lem-basis}
For each $P\in\{A,R,Q\}$, the system consisting of (i) the families $K_{P}(V)$ at all nonempty $V$, (ii) identity families $\{\id_{V}\}$ at all $V$, and (iii) the null cover at $\varnothing$, is a pullback-stable coverage basis on $L$: pullback along any $U\subseteq V$ carries a basis family at $V$ to (a family generating the same sieve as) a basis family at $U$, and refining the members of a basis family by basis families yields a family whose generated sieve contains a basis-generated sieve.
\end{lemma}

\begin{proof}
We take the axioms in turn, with all cases.

\emph{Identity covers.} Clause (ii) supplies $\{\id_{V}\}$ at every context by fiat. This is the basis-level source of axiom (J1) below.

\emph{Pullback stability, main case.} Let $V\neq\varnothing$, $U\subseteq V$, $U\neq\varnothing$. By Proposition~\ref{ch07:prop-pullback-stability}, the pullback of the indexed family $K_{P}(V)$ along $U\to V$ is, block for block, exactly the indexed family $K_{P}(U)$: the member $V\cap B\to V$ pulls back to $(V\cap B)\cap U=U\cap B\to U$, using $U\subseteq V$. Chapter~7's Section 7.\ref{ch07:sec-pullback-stability} walked this identity through every context class---pullbacks to singletons, to each of the three doubleton kinds, to triples---so the stability holds not merely abstractly but with the pulled-back family computed and identified in each case.

\emph{Pullback stability, empty subcontext.} If $U=\varnothing$, each member of $K_{P}(V)$ pulls back to $\varnothing\cap B=\varnothing$, i.e., to $\id_{\varnothing}$, giving a family of identity arrows on $\varnothing$; independently, clause (iii) covers $\varnothing$ by the null family. Either way $\varnothing$ is covered, and the generated-sieve comparison below is unaffected.

\emph{Pullback of identity and null covers.} An identity family pulls back to an identity family. The null cover sits at the bottom object, along which the only pullback is via $\id_{\varnothing}$, fixing it.

\emph{Composition, same partition.} Refine each member $V\cap B\to V$ of $K_{P}(V)$ by its own family $K_{P}(V\cap B)$ (null cover if the piece is empty). By Lemma~\ref{ch07:lem-same-refinement}, block disjointness collapses the composite: $V\cap B\cap B'$ is $V\cap B$ when $B'=B$ and $\varnothing$ otherwise, so the composite family is the original together with redundant empty pieces. An empty piece $\varnothing\to V$ contributes to a generated sieve only the arrows factoring through $\varnothing$---that is, the single arrow $\varnothing\to V$ itself---and the sieve generated by $K_{P}(V)$ already contains $\varnothing\to V$, since $\varnothing$ factors through any member (or, when all members are empty, is a member). Hence the composite generates \emph{the same} sieve as the original family, which is stability under composition in the strongest form.

\emph{Composition, with identity and null covers.} Refining by identity families changes nothing. Refining an empty piece by the null cover contributes no arrows, and the generated sieve again contains everything it contained before.
\end{proof}

\begin{remark}\label{ch09:rem-basis-strength}
The lemma's proof delivered more than the axioms require: pulled-back families are equal (not merely cofinal) to the target families, and same-partition composites generate equal (not merely finer) sieves. This exactness is what the marked-cover recomputation checks of Chapter~8 rely on, and it is why the certificate can verify pullback stability as a family-level identity rather than a sieve-level inclusion.
\end{remark}

Before using the basis, we record why checking sheaf conditions on basis families suffices for the generated topology---the reduction invoked constantly below, made explicit once.

\begin{lemma}[Basis check suffices]\label{ch09:lem-basis-suffices}
Let $J$ be the topology generated by a coverage basis on $L$ as in Lemma~\ref{ch09:lem-basis}, and let $F$ be a presheaf. If the gluing map of $F$ is a bijection for every basis family (including the null cover), then $F$ is a $J$-sheaf: its gluing map is a bijection for every $J$-covering sieve.
\end{lemma}

\begin{proof}[Proof sketch]
This is the standard equivalence between the sheaf condition for a coverage and for its generated topology \cite{stacks00ZC, johnstone}; we indicate the shape of the argument on our finite poset. The covering sieves of $J$ are reachable from basis-generated sieves by three closure steps, and the sheaf condition propagates along each. (i) \emph{Generated sieves.} For the sieve $\langle K\rangle$ generated by a basis family $K$, matching families for $\langle K\rangle$ restrict to matching families for $K$, and conversely a matching family for $K$ extends uniquely to $\langle K\rangle$: an arrow $T\to V$ of $\langle K\rangle$ factors through some member $S_{k}$, and one defines $x_{T}=x_{S_{k}}|_{T}$, well-defined because two factorizations $T\subseteq S_{k}$, $T\subseteq S_{l}$ give $T\subseteq S_{k}\cap S_{l}$, where the family's compatibility applies; identity covers make this trivial and the null cover has no arrows to extend over. Hence bijectivity for $K$ gives bijectivity for $\langle K\rangle$. (ii) \emph{Supersets.} Let $\sigma\subseteq\sigma'$ be sieves on $V$ with $\sigma$ covering, and suppose the gluing map of $F$ is bijective for $\sigma$ and for the pullback $g^{*}\sigma$ at every $g\in\sigma'$. A matching family for $\sigma'$ restricts to one for $\sigma$ and so glues to a unique candidate section; the candidate restricts correctly on all of $\sigma'$ because the two sides agree after pullback along the covering sieve $g^{*}\sigma$ of each $g\in\sigma'$, where gluing is injective. (iii) \emph{Local character.} Sieves admitted by (J3) are handled by the same two-step restriction-extension argument applied along the locally-covering data. On a finite poset every such propagation is a finite computation, and the exact-topology verifier of Chapter~11 checks the end result---sheafhood against every generating cover at all sixteen contexts, for each of $J_{A}$ and $J_{R}$---directly, so the reduction is corroborated as well as cited. (The $168$ covering sieves counted in Corollary~\ref{ch09:cor-census} belong to the join $J_{A}\vee J_{R}$, not to either factor.)
\end{proof}

\section{NT2-1: the generated topologies}\label{ch09:sec-nt21}

\begin{theorem}[NT2-1: transverse marked context coverages generate Grothendieck topologies]\label{ch09:thm-nt21}
The families $K_{A}$, with identity families and the null cover, generate a Grothendieck topology $J_{A}$ on $L$; likewise $K_{R}$ generates $J_{R}$, and $K_{Q}$ generates a topology $J_{\mathrm{atom}}$. Explicitly, for each of these bases the collection of sieves
\begin{multline*}
J_{P}(V)\;=\;\bigl\{\,\sigma \text{ a sieve on } V \;\bigm|\; \sigma \text{ contains the sieve generated by some}\\
\text{basis family obtained by iterated pullback and composition}\\
\text{of basis families at } V\,\bigr\}
\end{multline*}
satisfies the three Grothendieck axioms: (J1) the maximal sieve covers; (J2) covering sieves pull back to covering sieves; (J3) a sieve locally covered by a covering sieve covers.
\end{theorem}

\begin{proof}
This is the standard passage from a coverage basis to a topology \cite{stacks00ZC}, instantiated on Lemma~\ref{ch09:lem-basis}; we spell out how each axiom is discharged.

(J1) The identity family $\{\id_{V}\}$ is a basis family; the sieve it generates is the maximal sieve $m_{V}$ (every arrow into $V$ factors through $\id_{V}$), so $m_{V}\in J_{P}(V)$ for every $V$. At $\varnothing$, additionally, the null cover generates the empty sieve, which is also covering.

(J2) It suffices to check stability on basis-generated sieves and then close up. Let $\sigma$ be generated by a basis family at $V$ and let $U\subseteq V$. The pullback sieve $\iota^{*}\sigma$ along $\iota:U\to V$ consists of the subsets $T\subseteq U$ with $T\in\sigma$ (in the down-set reading of sieves on a poset). By Lemma~\ref{ch09:lem-basis}, the pullback of the generating family is the basis family at $U$ (or identities on $\varnothing$), and the sieve generated by the pulled-back family is contained in $\iota^{*}\sigma$: each pulled-back member $U\cap B$ satisfies $U\cap B\subseteq V\cap B\in\sigma$, hence $U\cap B\in\sigma$ by down-closure, hence lies in $\iota^{*}\sigma$. So $\iota^{*}\sigma$ contains a basis-generated covering sieve and is covering. For a general covering sieve $\sigma'\supseteq\sigma$, monotonicity of pullback gives $\iota^{*}\sigma'\supseteq\iota^{*}\sigma$, and supersets of covering sieves are covering by construction of $J_{P}$.

(J3) Suppose $\sigma\in J_{P}(V)$ and $\tau$ is a sieve on $V$ with $f^{*}\tau\in J_{P}(U)$ for every $f:U\to V$ in $\sigma$. We must show $\tau\in J_{P}(V)$. In the generated topology this is closure under cover composition: unwinding, $\tau$ contains, over each member $S$ of a basis family generating (a sieve inside) $\sigma$, a basis-generated covering sieve of $S$; the composite family---members of basis families over members of a basis family---generates a sieve contained in $\tau$; and Lemma~\ref{ch09:lem-basis}'s composition clause shows such composites are again reachable from the basis by the operations under which $J_{P}$ is by definition closed. Hence $\tau$ contains a covering sieve and is covering. (For the specific bases at hand the composites were computed exactly: same-partition composites regenerate the same sieve, by Lemma~\ref{ch07:lem-same-refinement}; this is the only composition shape needed for a single $J_{P}$.)
\end{proof}

\begin{remark}[Ledger record]\label{ch09:rem-ledger-nt21}
Theorem~\ref{ch09:thm-nt21} is registered as:
\begin{center}
\begin{tabular}{@{}ll@{}}
\toprule
Ledger id & \texttt{THM-NT2-MARKED-COVERAGE} (version 0.1) \\
Statement kind & \texttt{THEOREM} \\
Proof status & \texttt{INFORMAL\_PROOF\_PRESENT} \\
Review state & \texttt{IN\_REVIEW} \\
Migration disposition & \texttt{IN\_REVIEW} \\
Novelty status & \texttt{NOVELTY\_UNASSESSED} \\
Admission ring & 4 \\
\bottomrule
\end{tabular}
\end{center}
The registered statement also covers the provenance half of the construction: ``Parallel immutable mark records preserve cover kind, source authority, evidence references, and pullback ancestry even when the generated unmarked families coincide''---that half is Chapter~8's, and the record's claim boundary keeps the two halves apart, verbatim: \emph{``The topology theorem is mathematical; the hash-stable provenance behavior is an executable serialization invariant. Neither is evidence that either marking measures an empirical property.''}
\end{remark}

\begin{boundary}\label{ch09:boundary-nt21}
Per the record's claim boundary just quoted: nothing in Theorem~\ref{ch09:thm-nt21} or its proof interprets the ACBE or RPM labels empirically. The theorem is about two pullback-stable partition bases on a finite poset; the names of the partitions do no mathematical work.
\end{boundary}

\section{Propriety: the generators contain no identities}\label{ch09:sec-proper}

The legacy collapse (Chapter~5) turned on a single mechanism: a sieve containing an identity, or a split epimorphism composable to an identity, is maximal. The successor's claim to have escaped is therefore exactly the claim that its global generators are \emph{proper}, and it calls for a careful proof rather than an inspection.

\begin{proposition}[NT2 global covers are proper]\label{ch09:prop-proper}
At the global context, both pieces of $K_{A}(\Omega)=\{A_{0}\to\Omega,\ A_{1}\to\Omega\}$ and both pieces of $K_{R}(\Omega)=\{R_{0}\to\Omega,\ R_{1}\to\Omega\}$ are proper subcontexts of $\Omega$, and the sieves these families generate are nonmaximal.
\end{proposition}

\begin{proof}
\emph{Properness.} $|A_{i}|=|R_{j}|=2<4=|\Omega|$, so each piece is a proper subset (rows $A_{0}$, $A_{1}$, $R_{0}$, $R_{1}$ of Table 7.\ref{ch07:tab-contexts}).

\emph{Nonmaximality, by the factoring argument.} Let $\sigma$ be the sieve on $\Omega$ generated by $K_{A}(\Omega)$. By definition of generation, an arrow $U\to\Omega$ lies in $\sigma$ if and only if it factors through some member, i.e., $U\subseteq A_{0}$ or $U\subseteq A_{1}$. Suppose toward a contradiction that $\sigma$ were maximal; then $\id_{\Omega}\in\sigma$, so $\Omega\subseteq A_{0}$ or $\Omega\subseteq A_{1}$. But properness gives $A_{i}\subsetneq\Omega$, so neither containment holds: no proper subset of a set contains it. Hence $\id_{\Omega}\notin\sigma$ and $\sigma$ is nonmaximal. The identical argument applies to $K_{R}(\Omega)$ with $R_{j}$ in place of $A_{i}$. (Concretely, $\sigma$ misses more than the identity: it contains none of the arrows from $\Omega$, the four triples, the two RPM blocks, or the two diagonals, since none of those is contained in a single ACBE block---the RPM blocks and the diagonals each meet both ACBE blocks by Remark~\ref{ch07:rem-doubleton-split}, and anything of size $\ge 3$ exceeds a block's size. Its seven members are $\varnothing$, the four singletons, and $A_{0}$, $A_{1}$.)
\end{proof}

\begin{remark}[The exact point of escape]\label{ch09:rem-escape}
Compare mechanism against mechanism. The legacy ACBE clause contained the reflexive descent map $\id_{X}$; any sieve containing an identity is maximal, since every $f$ factors as $\id\circ f$. The legacy RPM clause contained the counit $\varepsilon_{X}$ of a section-retraction pair, and sieve closure under precomposition then contains $\varepsilon_{X}\circ\eta_{X}=\id_{X}$: maximal again. The successor's generators contain no identity (both pieces are proper) and no split epimorphism onto the target (in a poset, every arrow is monic, and an arrow $U\to\Omega$ admits a section only if $\Omega\subseteq U$, forcing $U=\Omega$). Both legacy failure mechanisms are thus structurally unavailable---and the executable certificate closes the loop by re-verifying the legacy mechanisms themselves as sanity checks (\texttt{legacy\_acbe\_identity\_forces\_maximal\_sieve}, \texttt{legacy\_rpm\_split\_epi\_forces\_identity}) alongside the propriety of the new generators.
\end{remark}

\begin{remark}[Ledger record]\label{ch09:rem-ledger-proper}
Proposition~\ref{ch09:prop-proper} is registered as:
\begin{center}
\begin{tabular}{@{}ll@{}}
\toprule
Ledger id & \texttt{THM-NT2-NONTRIVIAL-GLOBAL-COVERS} (version 0.1) \\
Statement kind & \texttt{THEOREM} \\
Proof status & \texttt{INFORMAL\_PROOF\_PRESENT} \\
Review state & \texttt{IN\_REVIEW} \\
Migration disposition & \texttt{IN\_REVIEW} \\
Novelty status & \texttt{NOVELTY\_UNASSESSED} \\
Admission ring & 4 \\
\bottomrule
\end{tabular}
\end{center}
Claim boundary, verbatim: \emph{``Nonmaximality repairs the inherited formal collapse; it does not establish semantic correctness of the labels.''}
\end{remark}

\section{NT2-2: subcanonicity of the context factor}\label{ch09:sec-nt22}

A topology is subcanonical when every representable presheaf is a sheaf \cite{stacks00WO}. On the poset $L$, the representable $yW=\Hom_{L}(-,W)$ has a transparent description: $yW(V)$ is a one-element set if $V\subseteq W$ and empty otherwise, with the (unique possible) restriction maps. Subcanonicity is the statement that these sixteen truth-valued presheaves already glue correctly along the new covers---a basic sanity property, since a site whose own representables fail its sheaf condition supports no Yoneda-style reasoning inside its sheaf category.

\begin{theorem}[NT2-2: both context topologies are subcanonical]\label{ch09:thm-nt22}
Every representable presheaf on $L$ is a $J_{A}$-sheaf and a $J_{R}$-sheaf. Hence $J_{A}$ and $J_{R}$ are subcanonical.
\end{theorem}

\begin{proof}
By the basis criterion it suffices to verify the sheaf condition for the generating families and the null cover \cite{stacks00ZC}; supersets and saturations add no obligations beyond those closed under by the topology construction. Fix $W\in L$ and write $y=yW$. Fix $P\in\{A,R\}$.

\emph{Null cover at $\varnothing$.} The matching-family set over the empty family is a singleton (there is exactly one empty family of sections, and its compatibility condition is vacuous). The sheaf condition thus requires $y(\varnothing)$ to be a singleton, and it is: $\varnothing\subseteq W$ always, so $y(\varnothing)$ has exactly one element.

\emph{Generating cover $K_{P}(V)$ at nonempty $V$.} A matching family assigns to each block $B\in P$ a section over the piece $V\cap B$, compatibly on overlaps. Since $y$ takes only empty or singleton values, both injectivity and the compatibility bookkeeping simplify, but we track them fully.

\emph{Injectivity of the gluing map.} $y(V)$ has at most one element, so any map out of it is injective.

\emph{Surjectivity.} Suppose a matching family exists; we must produce a section over $V$, i.e., show $V\subseteq W$. For each block $B$ with $V\cap B\neq\varnothing$: the family provides an element of $y(V\cap B)$, so $V\cap B\subseteq W$. For blocks with $V\cap B=\varnothing$: trivially $V\cap B=\varnothing\subseteq W$. Hence every piece is contained in $W$, and by joint exhaustiveness (Proposition~\ref{ch07:prop-exhaustive})
\[
V=\bigcup_{B\in P}(V\cap B)\subseteq W .
\]
So $y(V)$ is nonempty; its unique element restricts to the unique element of each $y(V\cap B)$, necessarily matching the given family since those sets are singletons.

\emph{Conversely}, if $y(V)$ is nonempty ($V\subseteq W$), its element restricts to a matching family, and we have shown every matching family arises this way, uniquely. The gluing map is therefore a bijection for every generator, for every representable, and both generated topologies are subcanonical.
\end{proof}

The certificate corroborates the theorem as \texttt{representables\_are\_sheaves}: all sixteen representables checked against every generating cover of both partitions.

\begin{remark}[Ledger record]\label{ch09:rem-ledger-nt22}
\begin{center}
\begin{tabular}{@{}ll@{}}
\toprule
Ledger id & \texttt{THM-NT2-SUBCANONICAL-CONTEXT-FACTOR} (version 0.1) \\
Statement kind & \texttt{THEOREM} \\
Proof status & \texttt{INFORMAL\_PROOF\_PRESENT} \\
Review state & \texttt{IN\_REVIEW} \\
Migration disposition & \texttt{IN\_REVIEW} \\
Novelty status & \texttt{NOVELTY\_UNASSESSED} \\
Admission ring & 4 \\
\bottomrule
\end{tabular}
\end{center}
Claim boundary, verbatim: \emph{``This applies to the context factor L only. It does not make the inherited jointly-surjective topology on Noetica subcanonical.''}
\end{remark}

\begin{boundary}[What NT2-2 does not repair]\label{ch09:boundary-nt22}
We expand the boundary beyond the ledger sentence, because the temptation it blocks is real. Chapter~6 exhibited a model of the registered legacy axioms in which the jointly-surjective topology on $\cC$ is \emph{not} subcanonical: subcanonicity of $(\cC,J_{\mathrm{legacy}})$ is not derivable. Theorem~\ref{ch09:thm-nt22} concerns the context factor $L$ alone; the two results coexist without tension, and the second neither repairs nor bears on the first. The operative consequence for the whole program: a distinguished legacy object $X$ is represented in the legacy sheaf topos by its sheafification $a_{\mathrm{legacy}}(yX)$, never by silently treating $yX$ as a sheaf \cite{stacks00WO}. No amount of good behavior on the finite factor changes that rule.
\end{boundary}

\section{NT2-3: distinct, non-identity modalities}\label{ch09:sec-nt23}

The heart of the finite theory: the two topologies are not merely well-defined and subcanonical but \emph{different}, and different in the strong sense that their sheafification modalities are distinct and neither is trivial. Distinctness is witnessed, not deduced: we construct an explicit presheaf that satisfies the ACBE gluing condition and fails the RPM one, and symmetrically.

\subsection{The witness presheaves}

\begin{definition}[Witness presheaves]\label{ch09:def-witness}
Recall from Definition~\ref{ch07:def-complete} the complete-block index sets and their counts $\alpha(V)=\#\{i\mid A_{i}\subseteq V\}$, $\rho(V)=\#\{j\mid R_{j}\subseteq V\}$. Define $F_{A}:L^{\mathrm{op}}\to\Set$ by
\[
F_{A}(V)\;=\;\bigl\{\,\text{functions } b:\{\,i\in\{0,1\}\mid A_{i}\subseteq V\,\}\to\{0,1\}\,\bigr\},
\]
the set of \emph{bit assignments to the complete ACBE blocks contained in $V$}. For $U\subseteq V$, the restriction map $F_{A}(V)\to F_{A}(U)$ sends $b$ to its restriction along the inclusion of index sets $\{i\mid A_{i}\subseteq U\}\subseteq\{i\mid A_{i}\subseteq V\}$ (Corollary~\ref{ch07:cor-blocks-monotone}): it \emph{forgets} the bits of blocks no longer completely contained. Define $F_{R}$ symmetrically, with RPM blocks in place of ACBE blocks; equivalently, $F_{R}$ is the transport of $F_{A}$ along the swap automorphism $S$ of Proposition~\ref{ch07:prop-swap}.
\end{definition}

By Table 7.\ref{ch07:tab-blocks}, the cardinalities are
\[
|F_{A}(V)|=2^{\alpha(V)}:\qquad
1 \text{ on } \varnothing,\ \text{singletons},\ R\text{-blocks},\ \text{diagonals};\quad
2 \text{ on } A\text{-blocks and triples};\quad
4 \text{ on } \Omega,
\]
and dually $|F_{R}(V)|=2^{\rho(V)}$.

\subsection{Functoriality first}

\begin{lemma}[$F_{A}$ and $F_{R}$ are presheaves]\label{ch09:lem-functorial}
The restriction maps of Definition~\ref{ch09:def-witness} satisfy the identity and composition laws, so $F_{A}$ and $F_{R}$ are functors $L^{\mathrm{op}}\to\Set$.
\end{lemma}

\begin{proof}
Identity: restricting along $V\subseteq V$ restricts $b$ along the identity inclusion of index sets, giving $b$. Composition: for $T\subseteq U\subseteq V$, the index sets satisfy $\{i\mid A_{i}\subseteq T\}\subseteq\{i\mid A_{i}\subseteq U\}\subseteq\{i\mid A_{i}\subseteq V\}$, and restriction of functions along a composite inclusion equals the composite of restrictions. The same argument applies to $F_{R}$.
\end{proof}

\begin{remark}[The functoriality-checked-first discipline]\label{ch09:rem-functorial-first}
The lemma looks unmissable, and that is the danger: a \emph{proposed} witness that assigns section sets per context but garbles restriction maps can appear to ``fail gluing'' for reasons that are actually failures to be a presheaf at all, silently invalidating a distinctness argument. The finite implementation therefore validates functoriality of any proposed witness across all $81$ restriction maps and all $256$ composable pairs $T\subseteq U\subseteq V$---the two counts of Theorem~\ref{ch07:thm-81} and Proposition~\ref{ch07:prop-chains} respectively---\emph{before} evaluating any sheaf condition, so that a nonfunctorial assignment can never pass as a witness. The certificate registers \texttt{acbe\_witness\_functorial} and \texttt{rpm\_witness\_functorial} as separate, prior checks to the four sheaf verdicts. The monograph observes the same discipline: Lemma~\ref{ch09:lem-functorial} precedes Theorem~\ref{ch09:thm-nt23}.
\end{remark}

\subsection{The theorem}

\begin{theorem}[NT2-3: distinct, non-identity modalities]\label{ch09:thm-nt23}
$F_{A}$ is a $J_{A}$-sheaf but not a $J_{R}$-sheaf, and $F_{R}$ is a $J_{R}$-sheaf but not a $J_{A}$-sheaf. Consequently:
\begin{enumerate}
\item the reflective subcategories $\Sh(L,J_{A})$ and $\Sh(L,J_{R})$ of $\PSh(L,\Set)$ are distinct;
\item the sheafification modalities $M_{A}$ and $M_{R}$ (sheafify, then include) are distinct;
\item neither $M_{A}$ nor $M_{R}$ is naturally isomorphic to the identity on $\PSh(L,\Set)$.
\end{enumerate}
\end{theorem}

\begin{proof}
We prove four assertions about the witnesses, then derive the consequences.

\emph{(a) $F_{A}$ satisfies the null-cover condition.} $F_{A}(\varnothing)$ is the set of functions from the empty index set to $\{0,1\}$, which is a singleton (the empty function), as the null cover requires.

\emph{(b) $F_{A}$ satisfies ACBE gluing at every nonempty $V$.} Fix the cover $K_{A}(V)=\{V\cap A_{0}\to V,\ V\cap A_{1}\to V\}$. First, localize the index sets. By Corollary~\ref{ch07:cor-blocks-monotone}, $A_{k}\subseteq V\cap A_{i}$ iff $k=i$ and $A_{i}\subseteq V$: each complete block of $V$ is complete in exactly one piece (its own), and a piece contains no other complete blocks. Hence the index set of $V$ is the disjoint union of the index sets of the two pieces:
\[
\{\,i\mid A_{i}\subseteq V\,\}\;=\;\{\,i\mid A_{i}\subseteq V\cap A_{0}\,\}\ \sqcup\ \{\,i\mid A_{i}\subseteq V\cap A_{1}\,\}.
\]
Second, dispose of overlaps. The two pieces intersect in $(V\cap A_{0})\cap(V\cap A_{1})=V\cap(A_{0}\cap A_{1})=\varnothing$, where $F_{A}(\varnothing)$ is a singleton by (a); the compatibility condition of a matching family over $K_{A}(V)$ is therefore automatic. A matching family is thus exactly a free pair: a bit assignment on the complete blocks of $V\cap A_{0}$ and one on those of $V\cap A_{1}$. By the disjoint-union decomposition, such pairs correspond exactly to bit assignments on the complete blocks of $V$; the correspondence is implemented by the gluing map (restrict a $b\in F_{A}(V)$ to each piece, i.e., forget the other piece's indices), which is therefore a bijection. This holds uniformly in $V$---including the degenerate cases: when a piece is empty, its section set is the singleton and its index set is empty, contributing nothing to either side; when $V$ has no complete blocks at all (singletons, RPM blocks, diagonals: Table 7.\ref{ch07:tab-blocks}), both sides are singletons.

\emph{(c) $F_{A}$ fails RPM gluing at $\Omega$: the four-versus-one computation.} Consider the cover $K_{R}(\Omega)=\{R_{0}\to\Omega,\ R_{1}\to\Omega\}$. Count both sides of the gluing map.

\emph{Global sections:} $\alpha(\Omega)=2$, so $F_{A}(\Omega)=\{0,1\}^{\{0,1\}}$ has exactly $2^{2}=4$ elements: the four bit-pairs $(b(0),b(1))\in\{00,01,10,11\}$.

\emph{Matching families:} each RPM block contains no complete ACBE block---$A_{i}\subseteq R_{j}$ would contradict transversality, since $A_{i}\cap R_{j}=\{\omega_{ij}\}$ is a singleton while $|A_{i}|=2$ (Corollary~\ref{ch07:cor-blocks-cross})---so $\alpha(R_{0})=\alpha(R_{1})=0$ and $F_{A}(R_{0})$, $F_{A}(R_{1})$ are singletons. The overlap is $R_{0}\cap R_{1}=\varnothing$ with $F_{A}(\varnothing)$ a singleton, so compatibility is automatic, and the matching-family set is a product of two singletons: exactly $1$ element.

\emph{The verdict:} the gluing map is a function from a four-element set to a one-element set. It is surjective but drastically non-injective: all four global sections restrict to the same (unique) matching family, because restriction to either RPM block forgets both bits. The sheaf condition requires a bijection; $4\neq 1$. Hence $F_{A}$ is not a $J_{R}$-sheaf.

\emph{(d) The symmetric facts for $F_{R}$.} Transport along the swap automorphism $S$ (Proposition~\ref{ch07:prop-swap}, Remark~\ref{ch07:rem-symmetry-use}): $S$ exchanges the partitions and carries $K_{A}$-covers to $K_{R}$-covers, and $F_{R}=F_{A}\circ S$ by construction, so (b) transports to ``$F_{R}$ satisfies RPM gluing everywhere'' and (c) transports to ``$F_{R}$ fails ACBE gluing at $\Omega$,'' with the same $4$-versus-$1$ count against the cover $K_{A}(\Omega)$.

\emph{Consequences.} (1) $F_{A}\in\Sh(L,J_{A})\setminus\Sh(L,J_{R})$ and $F_{R}\in\Sh(L,J_{R})\setminus\Sh(L,J_{A})$: each subcategory contains an object the other excludes, so they are distinct (indeed incomparable). (2) Reflectors onto distinct reflective subcategories are distinct. (3) Suppose $M_{A}\cong\id$. Sheafhood is invariant under isomorphism of presheaves, and $M_{A}F$ is always a $J_{A}$-sheaf; applying the assumed isomorphism at $F=F_{R}$ would make $F_{R}$ isomorphic to the $J_{A}$-sheaf $M_{A}F_{R}$, hence itself a $J_{A}$-sheaf, contradicting (d). So $M_{A}\not\cong\id$; symmetrically $M_{R}\not\cong\id$. (More pointedly: a reflector fixes, up to its unit isomorphism, exactly the objects of its subcategory, and each witness lies outside the other's subcategory.)
\end{proof}

The certificate records the four verdicts exactly as proved: \texttt{acbe\_witness\_is\_acbe\_sheaf} true, \texttt{acbe\_witness\_is\_rpm\_sheaf} false, \texttt{rpm\_witness\_is\_rpm\_sheaf} true, \texttt{rpm\_witness\_is\_acbe\_sheaf} false.

\begin{remark}[Where else the RPM gluing of $F_{A}$ fails]\label{ch09:rem-other-failures}
The proof used the failure at $\Omega$, and one failure suffices. But the census of Table 7.\ref{ch07:tab-blocks} locates the full failure set, and it is instructive: against $K_{R}(A_{i})=\{\{\omega_{i0}\}\to A_{i},\ \{\omega_{i1}\}\to A_{i}\}$, the context $A_{i}$ has $|F_{A}(A_{i})|=2$ while both singleton pieces and the empty overlap carry singletons, so the gluing map is $2$-to-$1$: RPM gluing of $F_{A}$ already fails at both ACBE blocks. It also fails at every triple. At $T_{ij}$ one has $|F_{A}(T_{ij})|=2$, because $A_{1-i}$ is complete in $T_{ij}$, while the two RPM pieces $T_{ij}\cap R_{j}=\{\omega_{1-i,j}\}$ and $T_{ij}\cap R_{1-j}=R_{1-j}$ both carry singleton section sets and meet in $\varnothing$, so the matching-family set is a singleton and the gluing map is again $2$-to-$1$. Indeed, by transversality no RPM piece of any context contains a complete ACBE block---$V\cap R_{j}\subseteq R_{j}$ and $\alpha(R_{j})=0$, hence $\alpha(V\cap R_{j})=0$, by Corollary~\ref{ch07:cor-blocks-cross} together with monotonicity, Corollary~\ref{ch07:cor-blocks-monotone}; see also Appendix~C, Finding~\ref{appC:fnd:novaluecrossing}---so the RPM matching-family set of $F_{A}$ is a singleton at \emph{every} context: the gluing map is a bijection exactly where $\alpha(V)=0$---at $\varnothing$, the four singletons, the two RPM blocks, and the two diagonals---and fails injectivity exactly where $\alpha(V)\ge 1$. The witness is thus not a boundary-case curiosity: it separates the topologies at seven of the sixteen contexts---the two ACBE blocks, the four triples, and $\Omega$---with the global context giving the widest gap, $4$ versus $1$. Appendix~C tabulates the pattern context by context.
\end{remark}

\begin{remark}[Ledger record]\label{ch09:rem-ledger-nt23}
\begin{center}
\begin{tabular}{@{}ll@{}}
\toprule
Ledger id & \texttt{THM-NT2-DISTINCT-MODALITIES} (version 0.1) \\
Statement kind & \texttt{THEOREM} \\
Proof status & \texttt{INFORMAL\_PROOF\_PRESENT} \\
Review state & \texttt{IN\_REVIEW} \\
Migration disposition & \texttt{IN\_REVIEW} \\
Novelty status & \texttt{NOVELTY\_UNASSESSED} \\
Admission ring & 4 \\
\bottomrule
\end{tabular}
\end{center}
Claim boundary, verbatim: \emph{``No commutation or noncommutation claim is made in this version.''} (The commutation question is settled---affirmatively---by the composition theorem of Chapter~10, registered as a separate record at version 0.2; the NT2-3 record deliberately does not anticipate it.)
\end{remark}

\subsection{The ambient-lift caveat}\label{ch09:sec-ambient-caveat}

Faithful reproduction of the ledger requires reporting a discrepancy between the registered statement and what this monograph asserts. The ledger statement for \texttt{THM-NT2-DISTINCT-MODALITIES} ends with the sentence: \emph{``External product with a representable yC\_0 lifts the witnesses to the ambient successor topos PSh(C x L), so neither ambient modality is the identity.''} The publication layer has since identified that sentence as unproved as stated, and this monograph, following the registered correction, does \emph{not} assert it. The theorem proved above is scoped to $\PSh(L,\Set)$.

\begin{boundary}[The external-product witness is false in general]\label{ch09:boundary-ambient}
The proposed lift takes a witness $F$ on $L$ and a representable $y\cC_{0}$ on $\cC$ and forms the external product presheaf on $\cC\times L$, with value $\Hom_{\cC}(c,\cC_{0})\times F(V)$ at $(c,V)$. Two obstructions:
\begin{enumerate}
\item \emph{Failure at the null cover.} In any reasonable lifted topology, the null cover of the $L$ coordinate covers $(c,\varnothing)$, forcing sheaves to take singleton value there. But the external product's value at $(c,\varnothing)$ is $\Hom_{\cC}(c,\cC_{0})\times F(\varnothing)\cong\Hom_{\cC}(c,\cC_{0})$, which need not be a singleton for a general $\cC$ and $\cC_{0}$. The candidate witness can thus fail to be a sheaf for the topology on the side where it was supposed to \emph{pass}, collapsing the intended contrast.
\item \emph{Duplicated representable coordinates.} Even away from the null cover, gluing computations in the $\cC$ coordinate can be obstructed by the representable factor when covers interact with it nontrivially; the finite counting argument of Theorem~\ref{ch09:thm-nt23} does not survive the product unexamined.
\end{enumerate}
A terminal $\cC$-presheaf in place of $y\cC_{0}$ is a candidate repair---its value is a singleton everywhere, removing obstruction (1)---but it becomes a proof only after the lifted topology on $\cC\times L$ is actually defined, $\cC$ is assumed nonempty, and naturality is checked in both coordinates. None of that is part of v0.1. \emph{Ambient distinctness of the modalities therefore remains \textsc{open}}; the proved result is Theorem~\ref{ch09:thm-nt23} on $\PSh(L,\Set)$, and no artifact of this monograph asserts the ambient statement.
\end{boundary}

\section{NT2-4: the join is the singleton topology}\label{ch09:sec-nt24}

The two topologies interact through their join---the coarsest topology refining both. On the finite lattice the join is computed exactly, and it is exactly what transversality suggests: the topology of the four-singleton partition.

\subsection{The transverse refinement}

\begin{lemma}[Common refinement is the four-singleton cover]\label{ch09:lem-refinement}
At the global context, refining the cover $K_{A}(\Omega)$ by the covers $K_{R}(A_{i})$ of its pieces (in either order of partitions) yields the family whose members are the four singletons:
\[
\{\,A_{i}\cap R_{j}\to\Omega\,\}_{i,j}\;=\;\{\,\{\omega_{ij}\}\to\Omega\,\}_{i,j}\;=\;K_{Q}(\Omega).
\]
More generally, at any nonempty $V$ the cross-partition composite has nonempty members exactly the singletons of $V$ (with one empty piece per atom of $\Omega\setminus V$), and at $V=\varnothing$ the composite is the null cover.
\end{lemma}

\begin{proof}
The global case is Example~\ref{ch07:ex-composite-global}; the general case is Lemma~\ref{ch07:lem-cross-refinement}, both resting on the transversality identity $A_{i}\cap R_{j}=\{\omega_{ij}\}$ (Theorem~\ref{ch07:thm-transverse}).
\end{proof}

\subsection{The join computation}

Order topologies on $L$ by $J\le J'$ iff every $J$-covering sieve is $J'$-covering; the join $J_{A}\vee J_{R}$ is the smallest topology above both. Write $J_{\mathrm{atom}}$ for the topology generated by the singleton partition $Q$ with the null cover, as in Theorem~\ref{ch09:thm-nt21}.

\begin{lemma}[Join identification]\label{ch09:lem-join}
$J_{A}\vee J_{R}=J_{\mathrm{atom}}$.
\end{lemma}

\begin{proof}
Both inequalities are proved directly.

\emph{($\le$): $J_{\mathrm{atom}}$ contains both $J_{A}$ and $J_{R}$, hence their join.} Since the join is the smallest upper bound, it suffices to show $J_{A}\le J_{\mathrm{atom}}$ and $J_{R}\le J_{\mathrm{atom}}$, and by the basis criterion it suffices to show every generating $K_{A}$- or $K_{R}$-sieve is $J_{\mathrm{atom}}$-covering. Fix nonempty $V$ and consider the sieve $\sigma$ generated by $K_{A}(V)$. Every singleton $\{\omega\}\subseteq V$ lies inside exactly one ACBE piece: $\omega\in A_{i}$ for exactly one $i$, so $\{\omega\}\subseteq V\cap A_{i}$, and by down-closure $\{\omega\}\in\sigma$. Thus $\sigma$ contains \emph{every} singleton of $V$, hence contains the sieve generated by the singleton family $K_{Q}(V)$. Since $\sigma\supseteq\langle K_{Q}(V)\rangle$ and $\langle K_{Q}(V)\rangle$ is $J_{\mathrm{atom}}$-covering by construction, monotonicity of covering sieves (Lemma~\ref{ch03:lem:monotone}) gives $\sigma\in J_{\mathrm{atom}}(V)$. The same argument verbatim handles $K_{R}(V)$, and the null covers of all three bases agree at $\varnothing$. Hence $J_{A},J_{R}\le J_{\mathrm{atom}}$ and so $J_{A}\vee J_{R}\le J_{\mathrm{atom}}$.

\emph{($\ge$): any topology containing both partial topologies covers by singletons.} Let $J$ be any Grothendieck topology with $J\ge J_{A}$ and $J\ge J_{R}$; we show every $K_{Q}$-generated sieve is $J$-covering, whence $J_{\mathrm{atom}}\le J$ and in particular $J_{\mathrm{atom}}\le J_{A}\vee J_{R}$. Fix nonempty $V$. Cover $V$ by its ACBE pieces: the sieve $\langle K_{A}(V)\rangle$ is $J$-covering since $J\ge J_{A}$. Let $\tau=\langle K_{Q}(V)\rangle$ be the singleton-generated sieve on $V$; we verify the hypothesis of (J3) for $\tau$ against the covering sieve $\langle K_{A}(V)\rangle$. A generic generator of $\langle K_{A}(V)\rangle$ is a piece $V\cap A_{i}\to V$ (arrows factoring through pieces reduce to the pieces by (J2) and monotonicity). Pull $\tau$ back to $V\cap A_{i}$: the result contains every singleton of $V\cap A_{i}$ (a singleton of the piece is a singleton of $V$, hence in $\tau$; its pullback membership is immediate). If $V\cap A_{i}=\varnothing$, the pullback contains the null cover's content and the empty context is $J$-covered because $J\ge J_{A}$ includes the null cover. If $V\cap A_{i}\neq\varnothing$, the sieve on $V\cap A_{i}$ generated by its singletons contains the sieve generated by $K_{R}(V\cap A_{i})$'s refinement---concretely, by transversality (Lemma~\ref{ch09:lem-refinement} applied at $V\cap A_{i}$), the RPM pieces of $V\cap A_{i}$ are
\[
(V\cap A_{i})\cap R_{j}=V\cap\{\omega_{ij}\},
\]
which are exactly singletons of $V\cap A_{i}$ or empty. So the pullback of $\tau$ to the piece contains the $K_{R}$-generated covering sieve of the piece, hence is $J$-covering (supersets again). The hypothesis of (J3) holds for every arrow of the covering sieve $\langle K_{A}(V)\rangle$, so $\tau\in J(V)$. At $V=\varnothing$ the singleton basis prescribes the null cover, which is $J$-covering by assumption. Hence $J_{\mathrm{atom}}\le J$.

Together: $J_{A}\vee J_{R}=J_{\mathrm{atom}}$.
\end{proof}

\subsection{The theorem, and the union-topology reading}

\begin{theorem}[NT2-4: the join resolves all four atoms]\label{ch09:thm-nt24}
The common refinement of the ACBE and RPM global covers is the four-singleton cover, and
\[
J_{A}\vee J_{R}\;=\;J_{\mathrm{atom}},
\]
the topology generated by the singleton partition with the null cover. Moreover $J_{\mathrm{atom}}$ admits the following exact description: a sieve $\sigma$ on $V$ is covering if and only if it contains every singleton of $V$---equivalently, if and only if the union of its members is $V$. The joined context coverage is thus the ordinary union coverage on the finite powerset lattice.
\end{theorem}

\begin{proof}
The refinement computation is Lemma~\ref{ch09:lem-refinement}; the join identification is Lemma~\ref{ch09:lem-join}. For the exact description: if $\sigma$ contains every singleton of $V$ then $\sigma\supseteq\langle K_{Q}(V)\rangle$ is covering. Conversely let $\sigma\in J_{\mathrm{atom}}(V)$ and fix $\omega\in V$; pulling back along $\{\omega\}\to V$ (axiom (J2)) gives a covering sieve on $\{\omega\}$. Covering sieves on a singleton context contain the identity: the generating family $K_{Q}(\{\omega\})$ has the full piece $\{\omega\}$ among its members (and empties), so its sieve is maximal, and every covering sieve on $\{\omega\}$---being a superset of a sieve containing $\id$, or reached by (J3) from such---contains $\id_{\{\omega\}}$; the null cover does not apply at nonempty contexts. Membership of $\id_{\{\omega\}}$ in the pullback says exactly $\{\omega\}\in\sigma$. Hence $\sigma$ contains every singleton of $V$. Finally, for a down-closed $\sigma$, containing every singleton of $V$ is equivalent to $\bigcup\sigma=V$: one direction takes unions; the other picks, for $\omega\in V=\bigcup\sigma$, a member $S\ni\omega$ and applies down-closure to $\{\omega\}\subseteq S$. At $V=\varnothing$ both the empty sieve (generated by the null cover) and the maximal sieve $\{\id_{\varnothing}\}$ cover, and both have union $\varnothing$, consistently.
\end{proof}

The certificate corroborates the refinement identity as \texttt{transverse\_refinement\_is\_singletons} and (in the marked layer) \texttt{joint\_refinement\_is\_four\_singletons}; the exact-topology verifier of Chapter~11 computes both presentations of the join and confirms their equality sieve-by-sieve.

\begin{corollary}[Covering-sieve census for the join]\label{ch09:cor-census}
For a context $V$ with $|V|=n$, the number of $J_{\mathrm{atom}}$-covering sieves on $V$ is: $2$ for $n=0$ (the empty and maximal sieves); and for $n\ge 1$, the number of down-closed families in $\mathcal{P}(V)$ containing all singletons, namely $1,2,9,114$ for $n=1,2,3,4$. The total over the sixteen contexts is
\[
1\cdot 2\;+\;4\cdot 1\;+\;6\cdot 2\;+\;4\cdot 9\;+\;1\cdot 114\;=\;2+4+12+36+114\;=\;168 .
\]
\end{corollary}

\begin{proof}
The characterization of covering sieves is Theorem~\ref{ch09:thm-nt24}. Counting: a down-closed family containing all singletons of $V$ automatically contains $\varnothing$ (for $n\ge1$) and is determined by its members of size $\ge 2$, which form a family closed downward within sizes $\ge 2$.

$n=1$: no sets of size $\ge 2$; exactly the maximal sieve: $1$.

$n=2$: the only candidate is the doubleton $V$ itself, in or out: $2$.

$n=3$: any subfamily of the three doubletons ($2^{3}=8$ choices), plus the triple $V$ only if all three doubletons are present ($+1$): $9$.

$n=4$: choose a set $X_{3}$ of triples, then doubletons $X_{2}$ containing every doubleton inside a chosen triple, then the top only if everything else is full. Summing over $X_{3}$ with free-doubleton counts as in the proof of Lemma~\ref{ch07:lem-downsets} (the same intersection patterns, one level up): $X_{3}=\varnothing$: $2^{6}=64$; $|X_{3}|=1$ ($4$ ways, forcing $3$ doubletons): $2^{3}=8$ each, $32$; $|X_{3}|=2$ ($6$ ways, forcing $5$): $2$ each, $12$; $|X_{3}|=3$ ($4$ ways, forcing all $6$): $1$ each, $4$; $|X_{3}|=4$: $1$. Subtotal $64+32+12+4+1=113$; the full sieve (top included) adds $1$: $114$.

The weighted total over the class sizes $1,4,6,4,1$ is the displayed sum, $168$.
\end{proof}

\begin{remark}[Reconciling the two 168s]\label{ch09:rem-168}
The exact-topology verifier reports $298$ sieves in total (derived from first principles in Theorem~\ref{ch07:thm-298}) of which $168$ are covering for the join---the number just derived. As flagged in Remark~\ref{ch07:rem-298}, this $168$ is numerically equal to, but conceptually unrelated to, the count $\delta(4)=168$ of \emph{all} sieves on the single object $\Omega$. The verifier computes the covering count twice, once from the join presentation $J_{A}\vee J_{R}$ and once from the singleton presentation $J_{\mathrm{atom}}$, and confirms both the equality of the two topologies and the count; the derivation above shows the number is forced.
\end{remark}

\begin{remark}[Ledger record]\label{ch09:rem-ledger-nt24}
\begin{center}
\begin{tabular}{@{}ll@{}}
\toprule
Ledger id & \texttt{THM-NT2-JOIN-DISCRETE-CONTEXT} (version 0.1) \\
Statement kind & \texttt{THEOREM} \\
Proof status & \texttt{INFORMAL\_PROOF\_PRESENT} \\
Review state & \texttt{IN\_REVIEW} \\
Migration disposition & \texttt{IN\_REVIEW} \\
Novelty status & \texttt{NOVELTY\_UNASSESSED} \\
Admission ring & 4 \\
\bottomrule
\end{tabular}
\end{center}
Claim boundary, verbatim: \emph{``This is a structural gluing result for Set-valued context data, not a probabilistic marginal-extension theorem.''} The boundary matters because the joined topology's ``ordinary gluing over the four atomic contexts'' invites a reading in terms of compatible marginals; Chapter~12's gluing-boundary discussion explains why that reading is blocked---pairwise-compatible probability marginals need not extend to a joint distribution, and no $\Set$-valued sheaf statement licenses a probabilistic consistency claim.
\end{remark}

\begin{boundary}[The \texttt{J\_atom} naming boundary]\label{ch09:boundary-atom}
The internal label $J_{\mathrm{atom}}$ (\texttt{J\_atom} in the registered artifacts) is a name for the topology generated by the four-singleton partition. It does \emph{not} assert that this topology is the standard category-theoretic \emph{atomic topology}---the topology on a category satisfying the right Ore condition in which every nonempty sieve covers \cite{maclane, johnstone}---and on $L$ it is not that: the sieve $\langle\varnothing\to\Omega\rangle=\{\varnothing\to\Omega\}$, i.e., the nonempty sieve generated by the empty piece, has union $\varnothing\neq\Omega$ and fails to cover, by Theorem~\ref{ch09:thm-nt24}. The preferred publication-layer name is $J_{\mathrm{singleton}}$; this monograph uses $J_{\mathrm{atom}}$ where continuity with the registered artifacts demands it, always with this boundary in force.
\end{boundary}

\section{Consolidated ledger view}\label{ch09:sec-ledger}

Table~\ref{ch09:tab-ledger} collects the status fields of every record in theorem ledger v0.1, including the three records whose mathematical content lives in other chapters: the resolved-identity theorem and the composition theorem (Chapter~10) and the backward-compatibility conjecture (Chapter~12). All statuses are reproduced verbatim from the ledger as of 2026-08-04; the monograph promotes none of them.

\begin{table}[htbp]
\centering
\caption{Theorem ledger v0.1: status fields, verbatim. All records carry admission ring 4. Abbreviations: kind T = \texttt{THEOREM}, C = \texttt{CONJECTURE}; proof IPP = \texttt{INFORMAL\_PROOF\_PRESENT}, PR = \texttt{PROOF\_REQUIRED}; review IR = \texttt{IN\_REVIEW}, UA = \texttt{UNASSESSED}; novelty NU = \texttt{NOVELTY\_UNASSESSED}.}
\label{ch09:tab-ledger}
\begin{tabular}{@{}llllllc@{}}
\toprule
Ledger id (\texttt{THM-NT2-\ldots}) & Ver. & Kind & Proof & Review & Novelty & Ch. \\
\midrule
\texttt{MARKED-COVERAGE} & 0.1 & T & IPP & IR & NU & 9 \\
\texttt{NONTRIVIAL-GLOBAL-COVERS} & 0.1 & T & IPP & IR & NU & 9 \\
\texttt{SUBCANONICAL-CONTEXT-FACTOR} & 0.1 & T & IPP & IR & NU & 9 \\
\texttt{DISTINCT-MODALITIES} & 0.1 & T & IPP & IR & NU & 9 \\
\texttt{RESOLVED-MODALITIES-IDENTITY} & 0.1 & T & IPP & IR & NU & 10 \\
\texttt{JOIN-DISCRETE-CONTEXT} & 0.1 & T & IPP & IR & NU & 9 \\
\texttt{PRODUCT-SITE-BACKWARD-COMPATIBILITY} & 0.1 & C & PR & UA & NU & 12 \\
\texttt{MODALITY-COMPOSITION} & 0.2 & T & IPP & IR & NU & 10 \\
\bottomrule
\end{tabular}
\end{table}

Three remarks on the table. First, the dependency structure recorded in the ledger runs exactly as this monograph's chapter order suggests: \texttt{MARKED-COVERAGE} has no dependencies; \texttt{NONTRIVIAL-GLOBAL-COVERS}, \texttt{SUBCANONICAL-CONTEXT-FACTOR}, and \texttt{JOIN-DISCRETE-CONTEXT} depend on it; \texttt{DISTINCT-MODALITIES} depends on the first two; \texttt{RESOLVED-MODALITIES-IDENTITY} and \texttt{MODALITY-COMPOSITION} depend on \texttt{DISTINCT-MODALITIES}; and the conjecture depends on \texttt{JOIN-DISCRETE-CONTEXT}. Second, the sole version-0.2 record is the composition theorem, whose v0.2 statement carries the commutation result proved in Chapter~10; its claim boundary scopes it to $\PSh(L,\Set)$ and lists the lift to $\PSh(\cC\times L)$, coherence in the $\cC$ coordinate, interaction with observational covers, proof-assistant verification, and independent review as \textsc{open}. Third, the sole conjecture---backward compatibility---carries proof status \texttt{PROOF\_REQUIRED} and review state \texttt{UNASSESSED}; its claim boundary states that backward compatibility ``remains a Gate 1 proof and implementation obligation, not a completed theorem,'' and that even if proved, the resolved context factor is four independent legacy channels without cross-channel semantics. Chapter~12 develops both.

\section{What the finite theorems establish, and what they do not}\label{ch09:sec-closing}

Assembled, the finite-context theory says: the transverse partition bases are genuine coverage bases (Lemma~\ref{ch09:lem-basis}) generating genuine Grothendieck topologies (Theorem~\ref{ch09:thm-nt21}) whose global covers are proper---escaping, by construction, both legacy collapse mechanisms (Proposition~\ref{ch09:prop-proper}); both topologies are subcanonical on the context factor (Theorem~\ref{ch09:thm-nt22}); they are distinct, with distinct non-identity sheafification modalities, witnessed by explicit finite presheaves whose gluing behavior is computed exactly---bijective ACBE gluing everywhere for $F_{A}$, a four-versus-one RPM failure at the global context, and the mirror facts for $F_{R}$ (Theorem~\ref{ch09:thm-nt23}); and their join is exactly the singleton topology, the ordinary union coverage on the finite powerset lattice, with $168$ covering sieves among the $298$ (Theorem~\ref{ch09:thm-nt24} and Corollary~\ref{ch09:cor-census}).

What they do not establish, in the ledger's own language and this chapter's boundaries: no empirical meaning for the ACBE or RPM marks (every record's claim boundary); no repair of legacy subcanonicity (Boundary~\ref{ch09:boundary-nt22}); no ambient distinctness on $\PSh(\cC\times L)$ (Boundary~\ref{ch09:boundary-ambient}); no commutation claim \emph{at this record version} (Remark~\ref{ch09:rem-ledger-nt23}); no standard-atomic-topology claim (Boundary~\ref{ch09:boundary-atom}); no probabilistic marginal-extension reading of the join (Remark~\ref{ch09:rem-ledger-nt24}); and no promotion of any status: every proof above is an \texttt{INFORMAL\_PROOF\_PRESENT} record's proof written at length, corroborated exhaustively by an executable certificate that is itself, per its own boundary, neither proof-assistant verification nor peer review nor a license for any claim beyond the mathematics on this one sixteen-object site.

Chapter~10 constructs the sheafification reflectors explicitly and proves the composition law $M_{A}M_{R}\cong M_{\mathrm{atom}}\cong M_{R}M_{A}$; Chapter~11 opens the certificate.
